\documentclass[12pt, twoside]{article}
%\documentclass[12pt, twoside]{article}
\usepackage[letterpaper, margin=1in, headsep=0.2in]{geometry}
\setlength{\headheight}{0.6in}
%\usepackage[english]{babel}
\usepackage[utf8]{inputenc}
\usepackage{microtype}
\usepackage{amsmath}
\usepackage{amssymb}
%\usepackage{amsfonts}
\usepackage[nomessages]{fp} %\FPeval{\var-name}{2*sin(pi/6)}
\usepackage{siunitx} %units in math. eg 20\milli\meter
\usepackage{yhmath} % for arcs, overparenth command
\usepackage{tikz} %graphics
\usetikzlibrary{quotes, angles, arrows, arrows.meta}
\usepackage{pgfplots}
\usepgfplotslibrary{fillbetween}
\pgfplotsset{compat=1.16}
\usepackage{graphicx} %consider setting \graphicspath{{images/}}
\usepackage{parskip} %no paragraph indent
\usepackage{enumitem}
\usepackage{multicol}
\usepackage{venndiagram}

\usepackage{fancyhdr}
\pagestyle{fancy}
\fancyhf{}
\renewcommand{\headrulewidth}{0pt} % disable the underline of the header
\raggedbottom
\hfuzz=2mm %suppresses overfull box warnings

\usepackage{hyperref}

\fancyhead[LE]{\thepage}
\fancyhead[RO]{Markscheme}
\fancyhead[LO]{La Scuola d'Italia / Dr.\ Huson / IB Math 11 \\* 28 October 2026}

\begin{document}
\subsubsection*{2.6 Quiz: Lines, Exponents and Radicals --- Markscheme}

\textbf{\large Version (a)} \hfill \emph{Total: 35 marks.}

\begin{enumerate}[itemsep=0.3cm]

%--- Problem 1: Gradient and equation of AB; perpendicular L through C; intersection F (not the midpoint) [10 marks] ---
\item[1.] [Maximum mark: 10]

\textbf{(a)} $m_{AB} = \dfrac{5 - 1}{4 - (-4)}$ \hfill \textbf{M1}

$m_{AB} = \dfrac{4}{8} = \dfrac{1}{2}$ \hfill \textbf{A1}

\emph{A bare answer earns M0A1. Wrong method: $\frac{\Delta x}{\Delta y} = 2$:
M0A0.}

\textbf{(b)} $y - 1 = \frac{1}{2}\left(x - (-4)\right)$ \hfill \textbf{M1}

$y - 1 = \frac{1}{2}(x + 4)$ \hfill \textbf{A1}

\emph{M1 for substituting their gradient and either point. Accept any
correct form, e.g.\ $y = \frac{1}{2}x + 3$ or $x - 2y = -6$; the answer
must be an equation. FT from (a).}

\textbf{(c)} $m_{L} = -2$ \hfill \textbf{A1}

$y - 6 = -2(x - 1)$ \hfill \textbf{M1A1}

\emph{A1 for the negative reciprocal of their gradient. M1 for substituting
their gradient and $C(1, 6)$; A1 for a correct equation in any form, e.g.\
$y = -2x + 8$. Wrong methods: the perpendicular bisector of $AB$,
$y = -2x + 3$, passes through the midpoint $(0,\ 3)$, not through $C$:
A1M0A0; $m_{L} = 2$ (no sign change) gives $y = 2x + 4$: A0M1A0. FT from
(a).}

\textbf{(d)} $\frac{1}{2}(x + 4) + 1 = -2(x - 1) + 6$ \hfill \textbf{M1}

$\frac{1}{2}x + 3 = -2x + 8$, \quad $\frac{5}{2}x = 5$

$x = 2$, \quad $y = 4$ \qquad $(2,\ 4)$ \hfill \textbf{A1A1}

\emph{M1 for solving their two lines together by any method (equating,
substitution, GDC). A1 for each coordinate. Wrong method: the midpoint of
$AB$, $(0,\ 3)$: M0A0A0. FT from (b) and (c).}

\newpage

%--- Problem 2: Cinema tickets: set up and solve a 2x2 system; number of children [7 marks] ---
\item[2.] [Maximum mark: 7]

\textbf{(a)} $2x + 3y = 47$ \hfill \textbf{A1}

$x + 4y = 41$ \hfill \textbf{A1}

\textbf{(b)} $x = 41 - 4y$, \quad $2(41 - 4y) + 3y = 47$ \hfill \textbf{(M1)}

$5y = 35$

adult ticket $x = 13$ euros \hfill \textbf{A1}

child ticket $y = 7$ euros \hfill \textbf{A1}

\emph{M1 for any valid method, including the GDC simultaneous-equation
solver; a correct GDC answer with no other working earns full marks. FT
from (a).}

\textbf{(c)} $2(13) + 7k = 61$ \hfill \textbf{M1}

$7k = 35$, \quad $k = 5$ children \hfill \textbf{A1}

\emph{Wrong method: swapping the prices, $2(7) + 13k = 61$, gives
$k = 3.62$, not a whole number: M0A0. FT from (b).}

\newpage

%--- Problem 3: Exponent laws; exact fractional negative power; same-base equations [10 marks] ---
\item[3.] [Maximum mark: 10]

\textbf{(a)} $\left(2a^{3}b\right)^{2} = 4a^{6}b^{2}$ \hfill \textbf{M1}

$4a^{6}b^{2} \times 3a^{-2}b^{-4} = 12a^{4}b^{-2}$ \hfill \textbf{A1}

$= \dfrac{12a^{4}}{b^{2}}$ \hfill \textbf{A1}

\emph{M1 for raising every factor of the bracket to the power $2$. The
final A1 needs positive exponents only. Wrong method: leaving the $2$
unsquared gives $\frac{6a^{4}}{b^{2}}$: M0A0A0.}

\textbf{(b)} $\left(\dfrac{8}{27}\right)^{-\frac{2}{3}} = \left(\dfrac{27}{8}\right)^{\frac{2}{3}} = \left(\dfrac{3}{2}\right)^{2}$ \hfill \textbf{M1}

$= \dfrac{9}{4}$ \hfill \textbf{A1}

\emph{Accept $2\frac{1}{4}$. $2.25$ from the GDC with no working earns
M0A1. Wrong method: ignoring the negative sign gives $\frac{4}{9}$:
M0A0.}

\textbf{(c)} $\left(3^{2}\right)^{x - 1} = \left(3^{3}\right)^{x - 3}$ \hfill \textbf{M1}

$2x - 2 = 3x - 9$ \hfill \textbf{A1}

$x = 7$ \hfill \textbf{A1}

\emph{M1 for writing both sides as powers of $3$. Wrong method:
$2x - 2 = 3x - 3$ (the $3$ not multiplied through) gives $x = 1$:
M1A0A0.}

\textbf{(d)} $2^{x + 3} = 2^{-5}$ \hfill \textbf{M1}

$x + 3 = -5$, \quad $x = -8$ \hfill \textbf{A1}

\emph{M1 for writing $\frac{1}{32}$ as $2^{-5}$. Wrong method:
$\frac{1}{32} = 2^{5}$ gives $x = 2$: M0A0.}

\newpage

%--- Problem 4: Surds: collect like surds, conjugate product, hence rationalize, linear equation with surd coefficients [8 marks] ---
\item[4.] [Maximum mark: 8]

\textbf{(a)} $5\sqrt{2} + 2\sqrt{2} - 4\sqrt{2}$ \hfill \textbf{M1}

$= 3\sqrt{2}$ \hfill \textbf{A1}

\emph{M1 for writing at least two of the surds as multiples of $\sqrt{2}$.
Wrong method: $\sqrt{50 + 8 - 32} = \sqrt{26}$: M0A0.}

\textbf{(b)} $9 - 3\sqrt{5} + 3\sqrt{5} - 5$ \hfill \textbf{M1}

$= 4$ \hfill \textbf{A1}

\emph{Accept $3^{2} - \left(\sqrt{5}\right)^{2}$ for the M1. Wrong method:
$9 + 5 = 14$: M0A0.}

\textbf{(c)} $\dfrac{8}{3 - \sqrt{5}} \times \dfrac{3 + \sqrt{5}}{3 + \sqrt{5}} = \dfrac{8\left(3 + \sqrt{5}\right)}{4}$ \hfill \textbf{M1}

$= 6 + 2\sqrt{5}$ \hfill \textbf{A1}

\emph{Accept $a = 6$, $b = 2$. FT from (b).}

\textbf{(d)} $4 = 3x\sqrt{2} - x\sqrt{2} = 2x\sqrt{2}$ \hfill \textbf{M1}

$x = \dfrac{4}{2\sqrt{2}} = \dfrac{2}{\sqrt{2}} = \sqrt{2}$ \hfill \textbf{A1}

\emph{M1 for collecting the $x$-terms on one side. The A1 needs
$\sqrt{2}$: $\frac{2}{\sqrt{2}}$ or $1.41$ earns A0.}

\end{enumerate}

\newpage

\textbf{\large Version (b)} \hfill \emph{Total: 38 marks.}

\begin{enumerate}[itemsep=0.3cm]

%--- Problem 1: Line 3x + 4y = 24 with intercepts A, B; perpendicular from P(-1, -3); foot F; area of triangle PAB [12 marks] ---
\item[1.] [Maximum mark: 12]

\textbf{(a)} $4y = -3x + 24$, \quad $y = -\frac{3}{4}x + 6$ \hfill \textbf{M1}

$m_{1} = -\dfrac{3}{4}$ \hfill \textbf{A1}

\emph{A bare answer earns M0A1.}

\textbf{(b)} $m_{2} = \dfrac{4}{3}$ \hfill \textbf{A1}

$y - (-3) = \frac{4}{3}\left(x - (-1)\right)$ \hfill \textbf{M1}

$y + 3 = \frac{4}{3}(x + 1)$ \hfill \textbf{A1}

\emph{Accept $y = \frac{4}{3}x - \frac{5}{3}$ or $4x - 3y = 5$. Wrong
methods: $m_{2} = -\frac{4}{3}$ (no sign change): A0M1A0; the perpendicular
bisector of $AB$, $y = \frac{4}{3}x - \frac{7}{3}$, passes through the
midpoint $(4,\ 3)$, not through $P$: A1M0A0. FT from (a).}

\textbf{(c)} $3x + 4\left(\frac{4}{3}(x + 1) - 3\right) = 24$ \hfill \textbf{M1}

$\frac{25}{3}x = \frac{92}{3}$, \quad $x = \dfrac{92}{25} = 3.68$ \hfill \textbf{A1}

$y = \dfrac{81}{25} = 3.24$ \hfill \textbf{A1}

\emph{M1 for solving the two equations together by any method, including
the GDC. Accept fractions or the exact decimals. Wrong method: the midpoint
of $AB$, $(4,\ 3)$: M0A0A0. FT from (b).}

\textbf{(d)} $A(8,\ 0)$, $B(0,\ 6)$, \quad $AB = \sqrt{8^{2} + 6^{2}} = 10$ \hfill \textbf{A1}

$PF = \sqrt{(3.68 + 1)^{2} + (3.24 + 3)^{2}}$ \hfill \textbf{M1}

$PF = \sqrt{60.84} = 7.8$ \hfill \textbf{A1}

$\text{Area} = \frac{1}{2}(10)(7.8) = 39$ \hfill \textbf{A1}

\emph{A1 for $AB = 10$, or for both intercepts if another method is used.
M1 for any valid area method that uses $P$: base $AB$ with height $PF$, an
enclosing rectangle ($81 - 13.5 - 24 - 4.5$), or the shoelace formula. A1
for the height $7.8$, or the correct intermediate of their method. Wrong
methods: the area of triangle $OAB$, $24$; $\frac{1}{2}(AB)(PA)$ with
$PA = 3\sqrt{10}$: M0. FT from (c).}

\newpage

%--- Problem 2: System of three linear equations in three unknowns; elimination with scaling (no equal coefficients) [7 marks] ---
\item[2.] [Maximum mark: 7]

Label the equations (1), (2), (3).

$3 \times (1) + 4 \times (2)$: \quad $6x + 9y - 12z + 20x - 8y + 12z = -60 + 136$ \hfill \textbf{M1}

$26x + y = 76$ \hfill \textbf{A1}

$(1) + 2 \times (3)$: \quad $2x + 3y - 4z + 6x + 8y + 4z = -20 + 22$ \hfill \textbf{M1}

$8x + 11y = 2$ \hfill \textbf{A1}

$y = 76 - 26x$, \quad $8x + 11(76 - 26x) = 2$, \quad $-278x = -834$

$x = 3$ \hfill \textbf{A1}

$y = 76 - 26(3) = -2$ \hfill \textbf{A1}

$3(3) + 4(-2) + 2z = 11$, \quad $z = 5$ \hfill \textbf{A1}

\emph{Any order of elimination is acceptable. First M1 for scaling and
combining two equations so that one variable cancels; A1 for the correct
two-variable equation. Second M1 for eliminating the same variable from a
different pair; A1 for the correct second equation. Then A1 for each of
$x = 3$, $y = -2$, $z = 5$, independently; the final values may come from
any valid method for the two-variable system. FT within the problem: a
student who makes one arithmetic slip in an elimination loses that A1 and
may earn the later A1s for values that follow correctly from their two
equations. Wrong methods: adding or subtracting equations without scaling,
so no variable cancels: M0; a sign error when subtracting (e.g.\ $-8y$
treated as $+8y$): M1A0. Answers from the GDC solver with no elimination
shown earn M0A0M0A0A1A1A1 (the question asks for elimination).}

\newpage

%--- Problem 3: Same-base equation; integers a, b from a surd product; equation quadratic in 3^x [11 marks] ---
\item[3.] [Maximum mark: 11]

\textbf{(a)} $\left(5^{2}\right)^{x + 1} = \left(5^{3}\right)^{2x - 3}$ \hfill \textbf{M1}

$2x + 2 = 6x - 9$ \hfill \textbf{A1}

$x = \dfrac{11}{4}$ \hfill \textbf{A1}

\emph{Accept $2.75$. M1 for writing both sides as powers of $5$. Wrong
methods: $2x + 1 = 6x - 9$ or $2x + 2 = 6x - 3$ (a bracket not multiplied
out) gives $\frac{5}{2}$ or $\frac{5}{4}$: M1A0A0.}

\textbf{(b)} $2a + 2b\sqrt{3} + a\sqrt{3} + 3b = 12 + 7\sqrt{3}$ \hfill \textbf{M1}

$2a + 3b = 12$, \quad $a + 2b = 7$ \hfill \textbf{A1}

$a = 3$, \quad $b = 2$ \hfill \textbf{A1}

\emph{Accept the conjugate route:
$\frac{12 + 7\sqrt{3}}{2 + \sqrt{3}} \times \frac{2 - \sqrt{3}}{2 - \sqrt{3}}
= \frac{24 - 12\sqrt{3} + 14\sqrt{3} - 21}{1} = 3 + 2\sqrt{3}$ (M1A1A1).
Accept trial with the expansion of $(2 + \sqrt{3})(3 + 2\sqrt{3})$ shown as
a check (M1A1A1). A bare answer earns M0A0A1.}

\textbf{(c)} $9^{x} = \left(3^{x}\right)^{2} = u^{2}$, so $3u^{2} - 82u + 27 = 0$ \hfill \textbf{M1}

$(3u - 1)(u - 27) = 0$

$u = \dfrac{1}{3}$ \hfill \textbf{A1}

$u = 27$ \hfill \textbf{A1}

$3^{x} = \dfrac{1}{3} \Rightarrow x = -1$ \hfill \textbf{A1}

$3^{x} = 27 \Rightarrow x = 3$ \hfill \textbf{A1}

\emph{M1 for the quadratic in $u$. The $u$ values may come from factorizing,
the formula or the GDC. Stopping at $u = \frac{1}{3}, 27$ earns
M1A1A1A0A0. Wrong method: reading $3^{x}$ as $3x$ gives $x = \frac{1}{9}, 9$:
A0A0 for those values.}

\newpage

%--- Problem 4: Arithmetic vs geometric: expressions, least n with v_n > u_n, least n with geometric sum > arithmetic sum [8 marks] ---
\item[4.] [Maximum mark: 8]

\textbf{(a)} $u_{n} = 7 + 3(n - 1)$ \hfill \textbf{A1}

$v_{n} = 2(1.5)^{n - 1}$ \hfill \textbf{A1}

\emph{Accept $u_{n} = 3n + 4$.}

\textbf{(b)} $2(1.5)^{n - 1} > 7 + 3(n - 1)$ \hfill \textbf{M1}

$u_{7} = 25$, $v_{7} = 22.8$; \quad $u_{8} = 28$, $v_{8} = 34.2$ \hfill \textbf{A1}

$n = 8$ \hfill \textbf{A1}

\emph{M1 for the inequality, or a GDC table or graph of both sequences. A1
for the values on both sides of the change. FT from (a): $v_{n} = 2(1.5)^{n}$
gives $n = 7$.}

\textbf{(c)} $\dfrac{2\left(1.5^{n} - 1\right)}{1.5 - 1} > \dfrac{n}{2}\left(2(7) + 3(n - 1)\right)$ \hfill \textbf{M1}

$n = 9$: $149.8 < 171$; \quad $n = 10$: $226.7 > 205$ \hfill \textbf{A1}

$n = 10$ \hfill \textbf{A1}

\emph{M1 for both sum formulas (or a GDC table of both sums). A1 for the
sums on both sides of the change. Wrong method: comparing terms instead of
sums gives $n = 8$: M0A0A0. FT from (a).}

\end{enumerate}
\end{document}
